% *** Authors should verify (and, if needed, correct) their LaTeX system  ***
% *** with the testflow diagnostic prior to trusting their LaTeX platform ***
% *** with production work. IEEE's font choices can trigger bugs that do  ***
% *** not appear when using other class files.                            ***
% The testflow support page is at:
% http://www.michaelshell.org/tex/testflow/


%%*************************************************************************
%% Legal Notice:
%% This code is offered as-is without any warranty either expressed or
%% implied; without even the implied warranty of MERCHANTABILITY or
%% FITNESS FOR A PARTICULAR PURPOSE!
%% User assumes all risk.
%% In no event shall IEEE or any contributor to this code be liable for
%% any damages or losses, including, but not limited to, incidental,
%% consequential, or any other damages, resulting from the use or misuse
%% of any information contained here.
%%
%% All comments are the opinions of their respective authors and are not
%% necessarily endorsed by the IEEE.
%%
%% This work is distributed under the LaTeX Project Public License (LPPL)
%% ( http://www.latex-project.org/ ) version 1.3, and may be freely used,
%% distributed and modified. A copy of the LPPL, version 1.3, is included
%% in the base LaTeX documentation of all distributions of LaTeX released
%% 2003/12/01 or later.
%% Retain all contribution notices and credits.
%% ** Modified files should be clearly indicated as such, including  **
%% ** renaming them and changing author support contact information. **
%%
%% File list of work: IEEEtran.cls, New_IEEEtran_how-to.pdf, bare_jrnl_new_sample4.tex,
%%*************************************************************************
\PassOptionsToPackage{unicode}{hyperref}
\PassOptionsToPackage{hyphens}{url}
\PassOptionsToPackage{dvipsnames,svgnames,x11names}{xcolor}
% Note that the a4paper option is mainly intended so that authors in
% countries using A4 can easily print to A4 and see how their papers will
% look in print - the typesetting of the document will not typically be
% affected with changes in paper size (but the bottom and side margins will).
% Use the testflow package mentioned above to verify correct handling of
% both paper sizes by the user's LaTeX system.
%
% Also note that the "draftcls" or "draftclsnofoot", not "draft", option
% should be used if it is desired that the figures are to be displayed in
% draft mode.
%
\documentclass[
  journal,
]{IEEEtran}%
% If IEEEtran.cls has not been installed into the LaTeX system files,
% manually specify the path to it like:
% \documentclass[journal]{../sty/IEEEtran}
\usepackage[cmex10]{amsmath}
\usepackage{amssymb}
\usepackage{iftex}
\ifPDFTeX
  \usepackage[T1]{fontenc}
  \usepackage[utf8]{inputenc}
  \usepackage{textcomp} % provide euro and other symbols
\else % if luatex or xetex
  \usepackage{unicode-math} % this also loads fontspec
  \defaultfontfeatures{Scale=MatchLowercase}
  \defaultfontfeatures[\rmfamily]{Ligatures=TeX,Scale=1}
\fi
%\usepackage{lmodern}
\ifPDFTeX\else
\fi
% Use upquote if available, for straight quotes in verbatim environments
\IfFileExists{upquote.sty}{\usepackage{upquote}}{}
\IfFileExists{microtype.sty}{% use microtype if available
  \usepackage[]{microtype}
  \UseMicrotypeSet[protrusion]{basicmath} % disable protrusion for tt fonts
}{}
\makeatletter
\parindent    1.0em
\ifCLASSOPTIONcompsoc
  \parindent    1.5em
\fi
\makeatother
\usepackage{xcolor}
\setlength{\emergencystretch}{3em} % prevent overfull lines

\setcounter{secnumdepth}{5}
% Make \paragraph and \subparagraph free-standing
\ifx\paragraph\undefined\else
  \let\oldparagraph\paragraph
  \renewcommand{\paragraph}[1]{\oldparagraph{#1}\mbox{}}
\fi
\ifx\subparagraph\undefined\else
  \let\oldsubparagraph\subparagraph
  \renewcommand{\subparagraph}[1]{\oldsubparagraph{#1}\mbox{}}
\fi

\usepackage{longtable,booktabs,array}
\newcounter{none} % for unnumbered tables
\usepackage{calc} % for calculating minipage widths
% Correct order of tables after \paragraph or \subparagraph
\usepackage{etoolbox}
\makeatletter
\patchcmd\longtable{\par}{\if@noskipsec\mbox{}\fi\par}{}{}
\makeatother
% Allow footnotes in longtable head/foot
\IfFileExists{footnotehyper.sty}{\usepackage{footnotehyper}}{\usepackage{footnote}}
\makesavenoteenv{longtable}
\usepackage{graphicx}
\makeatletter
\newsavebox\pandoc@box
\newcommand*\pandocbounded[1]{% scales image to fit in text height/width
  \sbox\pandoc@box{#1}%
  \Gscale@div\@tempa{\textheight}{\dimexpr\ht\pandoc@box+\dp\pandoc@box\relax}%
  \Gscale@div\@tempb{\linewidth}{\wd\pandoc@box}%
  \ifdim\@tempb\p@<\@tempa\p@\let\@tempa\@tempb\fi% select the smaller of both
  \ifdim\@tempa\p@<\p@\scalebox{\@tempa}{\usebox\pandoc@box}%
  \else\usebox{\pandoc@box}%
  \fi%
}
% Set default figure placement to htbp
\def\fps@figure{htbp}
\makeatother


% definitions for citeproc citations
\NewDocumentCommand\citeproctext{}{}
\NewDocumentCommand\citeproc{mm}{%
  \begingroup\def\citeproctext{#2}\cite{#1}\endgroup}
\makeatletter
 % allow citations to break across lines
 \let\@cite@ofmt\@firstofone
 % avoid brackets around text for \cite:
 \def\@biblabel#1{}
 \def\@cite#1#2{{#1\if@tempswa , #2\fi}}
\makeatother
\newlength{\cslhangindent}
\setlength{\cslhangindent}{1.5em}
\newlength{\csllabelwidth}
\setlength{\csllabelwidth}{3em}
\newenvironment{CSLReferences}[2] % #1 hanging-indent, #2 entry-spacing
 {\begin{list}{}{%
  \setlength{\itemindent}{0pt}
  \setlength{\leftmargin}{0pt}
  \setlength{\parsep}{0pt}
  % turn on hanging indent if param 1 is 1
  \ifodd #1
   \setlength{\leftmargin}{\cslhangindent}
   \setlength{\itemindent}{-1\cslhangindent}
  \fi
  % set entry spacing
  \setlength{\itemsep}{#2\baselineskip}}}
 {\end{list}}
\usepackage{calc}
\newcommand{\CSLBlock}[1]{\hfill\break\parbox[t]{\linewidth}{\strut\ignorespaces#1\strut}}
\newcommand{\CSLLeftMargin}[1]{\parbox[t]{\csllabelwidth}{\strut#1\strut}}
\newcommand{\CSLRightInline}[1]{\parbox[t]{\linewidth - \csllabelwidth}{\strut#1\strut}}
\newcommand{\CSLIndent}[1]{\hspace{\cslhangindent}#1}



\setlength{\emergencystretch}{3em} % prevent overfull lines

\providecommand{\tightlist}{%
  \setlength{\itemsep}{0pt}\setlength{\parskip}{0pt}}



 


\usepackage{physics}
\usepackage[version=3]{mhchem}
\usepackage{orcidlink}
\usepackage{float}
\floatplacement{table}{htb}
\makeatletter
\@ifpackageloaded{caption}{}{\usepackage{caption}}
\AtBeginDocument{%
\ifdefined\contentsname
  \renewcommand*\contentsname{Table of contents}
\else
  \newcommand\contentsname{Table of contents}
\fi
\ifdefined\listfigurename
  \renewcommand*\listfigurename{List of Figures}
\else
  \newcommand\listfigurename{List of Figures}
\fi
\ifdefined\listtablename
  \renewcommand*\listtablename{List of Tables}
\else
  \newcommand\listtablename{List of Tables}
\fi
\ifdefined\figurename
  \renewcommand*\figurename{Fig.}
\else
  \newcommand\figurename{Fig.}
\fi
\ifdefined\tablename
  \renewcommand*\tablename{Table}
\else
  \newcommand\tablename{Table}
\fi
}
\@ifpackageloaded{float}{}{\usepackage{float}}
\floatstyle{ruled}
\@ifundefined{c@chapter}{\newfloat{codelisting}{h}{lop}}{\newfloat{codelisting}{h}{lop}[chapter]}
\floatname{codelisting}{Listing}
\newcommand*\listoflistings{\listof{codelisting}{List of Listings}}
\makeatother
\makeatletter
\makeatother
\makeatletter
\@ifpackageloaded{caption}{}{\usepackage{caption}}
\@ifpackageloaded{subcaption}{}{\usepackage{subcaption}}
\makeatother
\usepackage[skip=2pt,font=footnotesize]{caption}
%\captionsetup{format=myformat}
\makeatletter
%\setlength{\cslhangindent}{0pt plus .5pt}
\providecommand{\bibfont}{\footnotesize}
\let\CSLReferences@rig=\CSLReferences
\renewcommand{\CSLReferences}[2]{
\bibfont\settowidth\csllabelwidth{[999]}
\CSLReferences@rig{#1}{#2}
\vskip 0.3\baselineskip plus 0.1\baselineskip minus 0.1\baselineskip%
}
\makeatother
\ifLuaTeX
  \usepackage{selnolig}  % disable illegal ligatures
\fi
\IfFileExists{bookmark.sty}{\usepackage{bookmark}}{\usepackage{hyperref}}
\IfFileExists{xurl.sty}{\usepackage{xurl}}{} % add URL line breaks if available
\urlstyle{same} % disable monospaced font for URLs
\hypersetup{
  pdftitle={A Sample Article Using quarto-ieee for IEEE Journal and Transactions},
  pdfauthor={David Folio; John Doe},
  pdfkeywords={IEEE, IEEEtran, journal, Quarto, Pandoc, template},
  colorlinks=true,
  linkcolor={blue},
  filecolor={Maroon},
  citecolor={Blue},
  urlcolor={Blue},
  pdfcreator={LaTeX via pandoc}}

% *** Do not adjust lengths that control margins, column widths, etc. ***
% *** Do not use packages that alter fonts (such as pslatex).         ***
% There should be no need to do such things with IEEEtran.cls V1.6 and later.
% (Unless specifically asked to do so by the journal or conference you plan
% to submit to, of course. )


% correct bad hyphenation here
\hyphenation{op-tical net-works semi-conduc-tor}

%
% paper title
% can use linebreaks \\ within to get better formatting as desired
% Do not put math or special symbols in the title.
% paper title
% can use linebreaks \\ within to get better formatting as desired
% Do not put math or special symbols in the title.
\title{A Sample Article Using \texttt{quarto-ieee} for IEEE Journal and
Transactions}

\author{
\thanks{The \texttt{quarto-ieee} template is freely available under the
MIT license on github: \url{https://github.com/dfolio/quarto-ieee}.}
David Folio\orcidlink{0000-0001-9430-6091},~\IEEEmembership{Member,
IEEE}
and~John Doe%
\thanks{David Folio is with Laboratoire Prisme, INSA Centre Val de
Loire, Bourges, 18800 France%
  Corresponding author: david.folio@insa-cvl.fr
}
\thanks{Unknown affiliation}
%by-author.affiliations
\thanks{John Doe is with Anonymous University%
}
%by-author.affiliations
\thanks{Template created June 23, 2023; revised
\texttt{r\ format(Sys.Date(),format=\textquotesingle{}\%B\ \%d,\ \%Y\textquotesingle{})}.}
}
\begin{document}

% The paper headers
\markboth{Journal XXX, Month Year}{D. Folio: A Sample Article Using
quarto-ieee}

% use for special paper notices

% make the title area
\maketitle

% As a general rule, do not put math, special symbols or citations
% in the abstract or keywords.
\begin{abstract}
This document describes the most common article elements and how to use
the \texttt{quarto-ieee} class with Pandoc/Quarto-Markdown to produce
files that are suitable for submission to IEEE journals.
\texttt{quarto-ieee} can produce conference, journal, and technical note
(correspondence) papers with a suitable choice of class options. It
intends to generate PDF and HTML outputs that closely mimick what IEEE
would generate.
\end{abstract}
% Note that keywords are not normally used for peerreview papers.
\begin{IEEEkeywords}
IEEE, IEEEtran, journal, Quarto, Pandoc, template
\end{IEEEkeywords}

% For peer review papers, you can put extra information on the cover
% page as needed:
% \ifCLASSOPTIONpeerreview
% \begin{center} \bfseries EDICS Category: 3-BBND \end{center}
% \fi
%
% For peerreview papers, this IEEEtran command inserts a page break and
% creates the second title. It will be ignored for other modes.
% \IEEEpeerreviewmaketitle


\section{Introduction}\label{sec-intro}

\IEEEPARstart{T}{his} file is intended to serve as a ``sample article
file'' for IEEE journal papers produced with (Pandoc/Quarto)-Markdown
using \texttt{IEEEtran.cls} version 1.8b and later for the PDF output.
It is based on \texttt{bare\_jrnl\_new\_sample4.tex} provided by IEEE
Publication Technology, Staff and available from
\url{https://template-selector.ieee.org/}. The most common elements are
covered in the simplified and updated instructions in
\texttt{New\_IEEEtran\_how-to.pdf}. For less common elements you can
refer back to the original \texttt{IEEEtran\_HOWTO.pdf}. It is assumed
that the reader has a basic working knowledge of {\LaTeX}
\citeproc{ref-mittelbach2023latex}{{[}1{]}} and of
(Pandoc/Quarto)-Markdown \citeproc{ref-MacFarlane_Pandoc}{{[}2{]}},
\citeproc{ref-Allaire_Quarto_2022}{{[}3{]}} markup.

\section{The Design, Intent, and Limitations of this
Templates}\label{the-design-intent-and-limitations-of-this-templates}

The \texttt{quarto-ieee} template is intended to \textbf{approximate the
final look and page length of the articles/papers} either in PDF output
or HTML output. \textbf{They are NOT intended to be the final produced
work that is displayed in print or on IEEEXplore\textsuperscript{®}}.
They will help to give the authors an approximation of the number of
pages and layout that will be in the final version.

\subsection{Unsuported feature and
limitations}\label{unsuported-feature-and-limitations}

Although most of the {\LaTeX} and \texttt{IEEEtran.cls} commands and
environment are supported, there are some limitations when trying to
export to a format other than PDF (e.g.~HTML output). For PDF output,
the reader can use the {\LaTeX} command directly. However, this may
break other output formats.\\
It can be can reported the following limitations of the
\texttt{quarto-ieee} template: - Several authors with same affiliation
produce weird output. In such case, it is recommended to use
\texttt{note} and \texttt{tex-author-no-affiliation:\ true}. - For
\texttt{PDF} output - \texttt{quarto-ieee} use a hack to handle the
\texttt{longtable} issue with 2-column {\LaTeX}
documents{[}\^{}issues-1023{]}. But, in some cases, a page overflow may
occur (see also \textbf{?@sec-tables}). - For \texttt{HTML} output - The
default Quarto toc is used, so the table of contents (toc) display is
not the same as on {[}IEEEXplore®{]}. - Footnote are put at the end of
document, while on {[}IEEEXplore®{]} there are placed in the accordion.
- Figures are not placed in the accordion. - {[}IEEEXplore®{]} specifics
(e.g.~citation metrics, etc.) - The \texttt{HTML} output is a Quarto
citeable article \citeproc{ref-quarto-citation}{{[}4{]}}, so a citation
appendix is automatically added to the article end.

\protect\phantomsection\label{refs}
\begin{CSLReferences}{0}{0}
\bibitem[\citeproctext]{ref-mittelbach2023latex}
\CSLLeftMargin{{[}1{]} }%
\CSLRightInline{F. Mittelbach and U. Fischer, \emph{The {LaTeX}
companion}, 3rd ed. {Addison Wesley Professional}, 2023. }

\bibitem[\citeproctext]{ref-MacFarlane_Pandoc}
\CSLLeftMargin{{[}2{]} }%
\CSLRightInline{J. MacFarlane, A. Krewinkel, and J. Rosenthal,
{``{Pandoc}.''} {[}Online{]}. Available:
\url{https://github.com/jgm/pandoc}}

\bibitem[\citeproctext]{ref-Allaire_Quarto_2022}
\CSLLeftMargin{{[}3{]} }%
\CSLRightInline{J. J. Allaire, C. Teague, C. Scheidegger, Y. Xie, and C.
Dervieux, {``{Quarto}.''} Jan-2022 {[}Online{]}. Available:
\url{https://github.com/quarto-dev/quarto-cli}}

\bibitem[\citeproctext]{ref-quarto-citation}
\CSLLeftMargin{{[}4{]} }%
\CSLRightInline{{``Quarto - {Creating Citeable Articles}.''}
{[}Online{]}. Available:
\url{https://quarto.org/docs/authoring/create-citeable-articles.html}.
{[}Accessed: 25-Oct-2023{]}}

\end{CSLReferences}


% Can use something like this to put references on a page
% by themselves when using endfloat and the captionsoff option.
\ifCLASSOPTIONcaptionsoff
  \newpage
\fi

% trigger a \newpage just before the given reference
% number - used to balance the columns on the last page
% adjust value as needed - may need to be readjusted if
% the document is modified later
%\IEEEtriggeratref{8}
% The "triggered" command can be changed if desired:
%\IEEEtriggercmd{\enlargethispage{-5in}}

% Uncomment when use biblatex with style=ieee
%\renewcommand{\bibfont}{\footnotesize} % for IEEE bibfont size

\pagebreak[3]
\begin{IEEEbiography}[\includegraphics{david-folio.png}]{David Folio}
Use \texttt{IEEEbiography} with figure as option and the author name as
the argument followed by the biography text.
\end{IEEEbiography}
\begin{IEEEbiographynophoto}{John Doe}
Use \texttt{IEEEbiographynophoto} and the author name as the argument
followed by the biography text.
\end{IEEEbiographynophoto}
% that's all folks
\end{document}

